\section{Discussion and Threats to Validity}
\label{sec:discussion}

\subsection{Underlying Mechanisms}
Our findings indicate that AST structural chunking does not inherently surpass sliding windows in requirement evidence retrieval. 
Two mechanisms explain this: 
First, \textit{information density}: 50-line windows (mean 41.0 LOC) capture declarations, imports, and helper calls, providing richer co-occurrence signals for neural encoders, whereas fine-grained AST chunks (mean 19.8 LOC) strip this lexical context. 
Second, \textit{contextual compensation}: prepending hierarchical headers injects the missing architectural namespace, allowing cross-encoders to match high-level requirement keywords against file and class names. 
Consequently, headers produce dramatic gains for AST chunks ($\Delta = +0.146$), but offer diminishing returns for line windows ($\Delta = +0.041$), which already encapsulate surrounding context.

\subsection{Threats to Validity}
\textbf{Internal Validity:} Mitigated by $100\%$ candidate pooling at top-5 across all 12 systems (100% Judged@5), 100% blinded dual annotation ($\kappa_w = 0.9822$, 98.07% exact agreement), standardized adjudication, and 0.0% parser fallback. 
\textbf{Statistical Conclusion Validity:} $N=25$ queries provide 80% power to detect paired differences of $0.080$--$0.130$ NDCG@10 ($\alpha = 0.05$ to $0.0125$). Significant effects (Comps 1, 2, 3, 5) exceed these thresholds; non-significant Comp 4 denotes absence of evidence, not equivalence. 
\textbf{Construct Validity:} Requirements were extracted from job postings using \texttt{gemini-1.5-flash} to reflect authentic industry skill expectations, though prompt phrasing could introduce stylistic artifacts. Relevance was evaluated on an established 3-point ordinal scale. 
\textbf{External Validity:} Evaluated on Java Spring Boot and React/TypeScript with BGE-M3 and BGE-Reranker-Base. Caution is warranted when generalizing to dynamically typed languages (Python) or distributed architectures.
